\begin{frame}[label=q5c4]{5c.4 · DataFrame: FULL OUTER JOIN}
  \small
  \pregunta{¿Qué departamentos aparecen solo entre menores o solo entre adultos con edad conocida?}
  \entrada{Dos agregados por \texttt{departamento}: menores de 18 y adultos. Excluir rango de edad NO ESPECIFICA.}
  \flujo{Filtrar edades conocidas\\Dos agregados}{\texttt{FULL OUTER JOIN}\\Clave: departamento}{\texttt{coalesce} a cero\\Sumar y calcular porcentaje}{26 departamentos\\Ambos grupos presentes}
  \begin{columns}[T,onlytextwidth]
    \column{0.55\textwidth}
    \textbf{Código Scala original (fragmento)}
    \codigo{sql/05c_4_full_menores_adultos}{27}{31}
    \column{0.42\textwidth}
    \textbf{Salida}\\[5pt]
    {\footnotesize
\begin{tabular}{@{}lll@{}}
\toprule
Departamento & Menores & Adultos\\
\midrule
Lima Met. & 19\,751 & 19\,281\\
Junín & 4\,525 & 2\,688\\
Cusco & 4\,299 & 2\,762\\
\bottomrule
\end{tabular}
    }
  \end{columns}
  \respuesta{Los 26 departamentos tienen menores y adultos. En Lima el porcentaje es 50.60\% sobre edad conocida; 3c.1 usa el total con edades no especificadas y obtiene 50.02\%.}
  \origen{sql}{05c_4_full_menores_adultos}
\end{frame}
